\chapter{Conclusion}

\section{Goals and innovations}
We have developed procedures for both routing and conflict resolution, and described them in detail. We have solved the majority of the instances in the DISPLIB benchmark library using these procedures. The routing decisions are done using a top-down approach, where we prevent trains going in the opposite direction from using the same resources, which not only reduce the delay, but also improve the safety aspect of the affected schedules.

The conflict resolution is done using the \emph{path-and-cycle formulation}, which is not used by any of the winning solutions for the DISPLIB, furthering the research of the train dispatching problem (TDP). Its main motivation is the removal of big-M constrains that weaken the MILP relaxation and cause numerical instability. While we have not surpassed state-of-the-art in for this problem, our procedure provides a significant reduction in time to achieve a conflict free solution, at the cost of suboptimal objective values. This should be sufficient for true real-time scheduling of trains.

Both routing and conflict resolution rely on an efficient aggregation of the initial data to reduce the problem size and the use of an innovative \emph{link graph} structure. The aggregation section, while inspired by previous solutions, clearly describes each structure and how to build it, and can be used for any other solution or a problem related to TDP.

We introduce the link graph to keep track of resource dependencies and to allow us to get structural implication rules for our resource precedence decisions. The formulated conditions for these rules are much easier to use in practice then conditions described in the relevant literature. The link graph can be updated incrementally to allow for efficient augmentation based on the current routing decisions.

\section{Future work}
There are two main improvements to our current solution. First is the implementation of a rerouting procedure where the initially selected routes would be adjusted based on the current suboptimal solution. Second is the cut cleanup procedure for the master MILP problem in the conflict resolution, where our problem becomes bloated with cuts that might no longer be necessary, slowing down the optimization.

Another related area of research is the use of the introduced link graph for related problems, such as the blocking job-shop problem, where it could be used to reduce the number of required variables and constrains, similar to how we use it for conflict resolution.